\chapter{Lezione 9 - 03/11/2026}

\section{Introduzione e Contesto}

Questa lezione risulterà complessa da seguire, in particolar modo per coloro che non hanno presenziato alle ultime lezioni, per i quali la comprensione degli argomenti trattati potrebbe rivelarsi pressappoco impossibile. L'esposizione seguirà rigorosamente la trattazione presente nel testo di \textbf{Kippenhahn}, alla quale verranno aggiunte alcune interpretazioni personali per chiarirne i concetti. Tutti i dettagli fondamentali sono inclusi, ma non sarà possibile riprendere dall'inizio gli argomenti già affrontati.

\section{La Teoria della Lunghezza di Mescolamento (MLT)}

Nelle lezioni precedenti, il lavoro si è concentrato sull'implementazione delle equazioni per il trasporto di energia per convezione, basandosi sul quadro teorico della \textbf{Teoria della Lunghezza di Mescolamento} (\textit{Mixing Length Theory}). Questo impianto formale conduce essenzialmente a un sistema di \textbf{cinque equazioni}.

\subsection{Flussi di Energia e Definizioni dei Gradienti}

Osservando in particolare la terza di queste equazioni, si evince che il flusso trasportato per convezione corrisponde, in sostanza, all'eccesso di energia interna dell'elemento convettivo trasportato. Tale elemento si muove con una velocità $v$, e la quantità $\rho v$ rappresenta il flusso di massa (il contenuto di massa) di tale elemento convettivo.

Il gradiente termico reale ($\nabla$) è determinato dalla combinazione dei flussi di energia (radiativo e convettivo). Tuttavia, è fondamentale chiarire un equivoco comune:
\begin{itemize}
\item È \textbf{errato} affermare che il gradiente reale corrisponda semplicemente alla somma del gradiente radiativo ($\nabla_{rad}$) e dell'elemento convettivo.
\item Una tale somma definirebbe soltanto un \textit{gradiente ipotetico}, ovvero il gradiente che sarebbe necessario qualora si dovesse trasportare l'intero flusso di energia (sia radiativo che convettivo) esclusivamente in regime radiativo.
\end{itemize}
Il gradiente reale emerge, in realtà, come una combinazione complessa e non lineare di tutti questi fattori.

\subsection{Considerazioni Energetiche e Scale Temporali}

La velocità dell'elemento convettivo è stata derivata attraverso considerazioni energetiche. Nello specifico, si è calcolata l'energia cinetica derivante dal lavoro compiuto dalla forza di galleggiamento (\textbf{forza di Archimede} o \textit{buoyancy force}). Tale lavoro si converte in energia cinetica e, calcolando questa quantità per unità di massa, è possibile ricavare la velocità dell'elemento in risalita.

Un altro aspetto cruciale riguarda la quantità di energia che l'elemento convettivo \textit{perde} lungo il suo percorso. Questo fattore determina l'\textbf{efficienza del trasporto convettivo di energia}: maggiore è la quantità di energia irradiata (persa) dall'elemento, minore sarà l'efficacia della convezione.

In questo ambito, il ruolo delle scale temporali è fondamentale:
\begin{itemize}
\item L'elemento convettivo cede energia a un tasso che è l'analogo della scala temporale di \textbf{Kelvin-Helmholtz} (ovvero il tempo necessario all'elemento per irradiare la propria energia termica).
\item Se questa scala temporale di perdita di energia è \textbf{più breve} rispetto al tempo di \textit{turnover} convettivo (il tempo impiegato dall'elemento per completare il ciclo di salita e discesa), l'energia persa durante il tragitto risulta trascurabile.
\end{itemize}

Quando ci si trova in quest'ultimo regime, si ottiene di fatto una \textbf{condizione adiabatica}: l'elemento convettivo raccoglie energia alla base della zona convettiva, inizia a risalire senza subire perdite significative durante il tragitto, giunge in prossimità della sommità dove infine cede la sua energia, per poi ridiscendere.

\subsection{I Gradienti e la Lunghezza di Mescolamento}

Attraverso questo ragionamento, è stata calcolata la differenza tra il gradiente di temperatura dell'elemento ($\nabla_{e}$) e il gradiente adiabatico ($\nabla_{ad}$). Questa quantità ($\nabla_{e} - \nabla_{ad}$) descrive come la temperatura dell'elemento risponda a una variazione della pressione esterna rispetto al caso ideale adiabatico (nel qual caso la differenza sarebbe pari a zero).

Utilizzando un'argomentazione analoga a quella impiegata per il trasporto radiativo, si può esprimere il rapporto tra questa differenza ($\nabla_{e} - \nabla_{ad}$) e la differenza tra il gradiente reale e il gradiente dell'elemento ($\nabla - \nabla_{e}$). Questo rapporto scala in funzione di vari parametri fisici: temperatura, opacità, densità e, soprattutto, la \textbf{lunghezza di mescolamento} (\textit{mixing length}).

La lunghezza di mescolamento è un \textbf{parametro libero} della teoria. Dal punto di vista fisico, può essere immaginata come la distanza percorsa dall'elemento convettivo affinché riesca a trasportare una quantità sufficiente di energia. Essendo un parametro libero, esso viene tipicamente fissato tramite vincoli osservativi oppure attraverso simulazioni idrodinamiche e magnetoidrodinamiche.

La calibrazione più celebre è quella basata sulla riproduzione del raggio solare, che definisce il \textbf{parametro della lunghezza di mescolamento solare calibrato}, il cui valore si attesta generalmente intorno a $2$. Nelle equazioni della teoria, questa quantità ha le dimensioni fisiche di una lunghezza, ma all'atto pratico viene scalata utilizzando l'\textbf{altezza di scala della pressione} ($H_p$). Di conseguenza, la lunghezza di mescolamento viene solitamente espressa in unità di altezza di scala della pressione.

\section{Riduzione del Sistema e Risoluzione Analitica}

Nella fase successiva della lezione precedente, il sistema originario di cinque equazioni è stato ridotto a sole \textbf{due equazioni}, introducendo quantità adimensionali. A partire da queste due, si è giunti a un'equazione cubica facilmente risolvibile. Tuttavia, questo procedimento matematico ha comportato l'introduzione di ulteriori parametri, offuscando in parte il quadro fisico precedentemente descritto.

\subsection{Le Equazioni del Trasporto Convettivo}

Lavorando con l'equazione per i gradienti e l'equazione della velocità, è stata derivata una prima relazione (chiamata \textit{Equazione A}), che lega la differenza tra il gradiente dell'elemento e il gradiente adiabatico alla radice quadrata della differenza tra gradiente reale e gradiente dell'elemento, pesata per un parametro che chiameremo $U$:
\begin{equation}
(\nabla_{e} - \nabla_{ad}) \propto U \sqrt{\nabla - \nabla_{e}}
\end{equation}
Il parametro $U$ racchiude tutti i fattori termodinamici e strutturali visti in precedenza (temperatura, opacità, densità, altezza di scala della pressione, lunghezza di mescolamento). Esso rappresenta essenzialmente una misura dell'\textbf{efficienza della convezione}. Analizzando l'equazione, è evidente che se le variabili si combinano per produrre un valore di $U$ molto piccolo, le differenze tra i gradienti tendono a zero, portando la convezione nel regime adiabatico.

Combinando le restanti tre equazioni, si ottiene un'ulteriore relazione che coinvolge la differenza tra il gradiente reale e il gradiente dell'elemento ($\nabla - \nabla_{e}$). Questa equazione dimostra che maggiore è questa differenza, maggiore sarà l'energia cinetica sviluppata dall'elemento in risalita.

Questa stessa equazione permette di legare il termine ($\nabla - \nabla_{e}$) alla differenza tra il gradiente reale e il gradiente radiativo ($\nabla - \nabla_{rad}$).
\begin{itemize}
\item Se questa differenza si annulla, significa che non vi è convezione: il gradiente reale risulta uguale al gradiente radiativo ($\nabla = \nabla_{rad}$) e il sistema ritorna al puro trasporto radiativo di energia.
\item Questa condizione di limite radiativo si verifica anche quando il parametro $U$ assume valori molto grandi. Affinché le equazioni restino bilanciate (non potendo le quantità divergere all'infinito), i termini legati alla convezione devono tendere a zero.
\end{itemize}
Sebbene questo procedimento algebrico risulti laborioso, costituisce un metodo estremamente efficace per raggruppare i vari termini e stabilire analiticamente quando la convezione risulti efficiente, inefficiente, o quanto essa sia preponderante rispetto al trasporto radiativo.

\subsection{L'Equazione Cubica}

Le due equazioni fondamentali sopra descritte possono essere combinate per generare un'\textbf{equazione cubica} (\textit{Equazione C}), la cui risoluzione analitica semplifica notevolmente il problema. Per procedere, si introducono alcune variabili di comodo, in particolare la variabile $x$, che rappresenta la soluzione dell'equazione in funzione del parametro $U$ e della differenza tra il gradiente reale e il gradiente adiabatico ($\nabla - \nabla_{ad}$).

Nel formalismo dell'equazione cubica è stato introdotto il termine $x^2$, definito in relazione a $U^2$ e alla differenza dei gradienti. La soluzione dell'equazione cubica fornisce il valore di $x$, dal quale è possibile ricavare direttamente il gradiente reale $\nabla$. Essendo noto a priori il gradiente adiabatico $\nabla_{ad}$, dal gradiente reale si può derivare a cascata qualsiasi altra variabile termodinamica del sistema.

\section{Casi Limite e Interpretazione Fisica dei Parametri}

La stabilità della soluzione e l'impatto dei parametri liberi dipendono fortemente dal regime in cui si trova la stella:
\begin{itemize}
\item \textbf{Valori elevati di $U$}: indicano che si è vicini alle condizioni radiative (trasporto radiativo).
\item \textbf{Valori molto piccoli di $U$}: indicano che si è vicini alle condizioni adiabatiche (convezione estremamente efficiente).
\end{itemize}
In questi due estremi (regime puramente radiativo o regime puramente convettivo adiabatico), i tre parametri liberi della teoria giocano un ruolo trascurabile, non inficiando la robustezza del modello. Le difficoltà subentrano invece nelle situazioni intermedie, ovvero nelle condizioni di \textbf{stabilità marginale} durante il transito da un regime all'altro. In quest'area di transizione, le debolezze intrinseche della teoria emergono: l'assunzione iniziale sulla lunghezza di mescolamento diventa cruciale e la soluzione finale risulta \textbf{fortemente dipendente dal modello} adottato.

\subsection{Il Significato del Parametro $U$}

Per comprendere più a fondo l'interpretazione fisica della teoria, esploriamo i casi limite. Definiamo innanzitutto una quantità $W$ come la differenza tra il gradiente radiativo e quello adiabatico:
\begin{equation}
W = \nabla_{rad} - \nabla_{ad}
\end{equation}
Supponiamo, per semplicità di notazione, di raggruppare i termini in due costanti generali $K_{rad}$ e $K_{conv}$, in modo tale che il flusso radiativo e il flusso convettivo possano essere espressi come:
\begin{align}
F_{rad} &\propto K_{rad} \nabla \
F_{conv} &\propto K_{conv} (\nabla - \nabla_{e})^{3/2}
\end{align}
Esplicitando il contenuto di $K_{rad}$ e $K_{conv}$, è possibile dimostrare agevolmente che il parametro $U$ risulta proporzionale al rapporto tra queste due quantità ($K_{rad} / K_{conv}$).

Inoltre, compiendo alcuni passaggi algebrici laboriosi ma concettualmente semplici, si dimostra che $U$ equivale al rapporto tra due scale temporali fondamentali:
\begin{equation}
U \approx \frac{\tau_{\text{free fall}}}{\tau_{\text{thermal adj}}}
\end{equation}
Ovvero, $U$ è il rapporto tra il tempo di caduta libera (\textit{free fall timescale}, il tempo impiegato dall'elemento per attraversare la zona convettiva) e il tempo di aggiustamento termico (\textit{thermal adjustment timescale}, il tempo necessario affinché l'elemento si equilibri termicamente con l'ambiente circostante).
\begin{itemize}
\item Se il tempo di aggiustamento termico è breve, $U$ assume un valore elevato, e la convezione risulta inefficace.
\item Se il tempo di aggiustamento termico è lungo, $U$ assume un valore piccolo, e la convezione risulta pienamente efficiente.
\end{itemize}

\subsection{Il Parametro Gamma ($\Gamma$)}

Infine, introduciamo un ulteriore parametro adimensionale che chiameremo $\Gamma$. Esso è definito come il rapporto tra la radice quadrata della differenza tra gradiente reale ed elemento, e due volte il parametro $U$:
\begin{equation}
\Gamma = \frac{\sqrt{\nabla - \nabla_{e}}}{2U}
\end{equation}
Con semplici manipolazioni algebriche, è possibile dimostrare che $\Gamma$ può essere riscritto in forma di rapporto fra gradienti:
\begin{equation}
\Gamma = \frac{\nabla - \nabla_{e}}{\nabla_{e} - \nabla_{ad}}
\end{equation}
Il significato di questa espressione è profondamente fisico:
\begin{itemize}
\item Il \textbf{numeratore} ($\nabla - \nabla_{e}$) è proporzionale all'\textbf{energia trasportata} dall'elemento convettivo.
\item Il \textbf{denominatore} ($\nabla_{e} - \nabla_{ad}$) è proporzionale all'\textbf{energia persa per radiazione} dall'elemento durante il suo tragitto.
\end{itemize}
Pertanto, anche il parametro $\Gamma$, così come $U$, fornisce una misura rigorosa dell'efficienza del trasporto convettivo, ma invece di utilizzare un rapporto tra scale temporali, utilizza direttamente il rapporto tra l'energia utilmente trasportata e l'energia dissipata per irraggiamento.

Sia $U$ che $\Gamma$ permettono dunque di proiettare l'intera visione fisica e dinamica del problema in grandezze che discriminano univocamente una convezione efficiente da una inefficiente.

\section{Analisi dei Casi Limite}

Per comprendere appieno le implicazioni fisiche della teoria, è opportuno analizzare in dettaglio i casi limite, valutando il comportamento del sistema agli estremi dei valori assumibili dal parametro adimensionale $U$.

\subsection{Regime di Convezione Efficiente ($U \ll 1$)}

Cosa accade quando il parametro $U$ è molto piccolo? Di conseguenza, il parametro $\Gamma$ risulta essere molto grande.
Osservando la prima equazione derivata precedentemente, se $U \to 0$, il termine ad esso proporzionale tende a zero. Questo implica fisicamente che il gradiente dell'elemento convettivo tende al gradiente adiabatico:
\begin{equation}
\nabla_{e} \to \nabla_{ad}
\end{equation}
Questo è esattamente ciò che si verifica nel \textit{bulk} (la regione più interna e massiccia) della zona convettiva.

Considerando la seconda equazione, affinché i termini rimangano finiti per un $U$ molto piccolo, è necessario che il gradiente reale tenda al gradiente dell'elemento convettivo ($\nabla \to \nabla_{e}$). Ma, sapendo che $\nabla_{e}$ tende a sua volta a $\nabla_{ad}$, possiamo dedurre che in questo regime si ottiene un gradiente reale essenzialmente pari a quello adiabatico:
\begin{equation}
\nabla \to \nabla_{ad}
\end{equation}

\subsection{Regime Radiativo ($U \gg 1$)}

Consideriamo il caso opposto: $U$ assume valori molto grandi ($U \to \infty$), il che implica un parametro $\Gamma$ molto piccolo.
Analizzando la seconda equazione, affinché il risultato si mantenga finito nonostante la divergenza di $U$, è necessario che il gradiente reale approssimi il gradiente radiativo:
\begin{equation}
\nabla \to \nabla_{rad}
\end{equation}
Di conseguenza, il flusso di energia tende a coincidere interamente con il \textbf{flusso radiativo}. Questo scenario descrive ciò che avviene alla sommità della zona convettiva solare, immediatamente al di sotto della fotosfera, dove il trasporto convettivo perde di efficacia.

\section{Condizione di Stabilità Marginale ($\Gamma = 1$)}

Un caso intermedio di fondamentale importanza è quello della \textbf{stabilità marginale}, che si verifica per $\Gamma = 1$. Questo valore segna il confine critico tra i diversi regimi di trasporto dell'energia.
Ponendo $\Gamma = 1$, si ottiene l'uguaglianza:
\begin{equation}
\sqrt{\nabla - \nabla_{e}} = 2U
\end{equation}
Con semplici passaggi algebrici, è possibile dimostrare che questa condizione corrisponde a una relazione lineare nel piano logaritmico, espressa da:
\begin{equation}
\log W = \log U + 1.23
\end{equation}
dove ricordiamo che $W = \nabla_{rad} - \nabla_{ad}$.
Questa equazione definisce una retta nel piano delle variabili.
\begin{itemize}
\item Un \textbf{aumento di $\Gamma$} comporta una traslazione di questa retta in una direzione del piano.
\item Una \textbf{diminuzione di $\Gamma$} comporta una traslazione nella direzione opposta.
\end{itemize}
In maniera analoga, esiste un'altra retta che interseca lo zero, corrispondente alla condizione $U^2 = x$, sebbene i dettagli matematici di questa specifica intersezione non vengano approfonditi in questa sede.

\section{L'Interazione dei Parametri $W$ e $U$}

Al di fuori dei casi limite, la fisica del trasporto è dominata dalla combinazione dei parametri $W$ e $U$:
\begin{itemize}
\item Il parametro \textbf{$W$} fornisce una misura dell'energia intrinseca dell'elemento convettivo in risalita.
\item Il parametro \textbf{$U$} quantifica il rapporto tra l'energia utilmente trasportata e l'energia persa (irradiata).
\end{itemize}

La classificazione dell'efficienza convettiva si riduce quindi a queste casistiche:
\begin{enumerate}
\item \textbf{Convezione molto efficiente:} L'elemento convettivo possiede molta energia e subisce perdite radiative nulle o trascurabili. In questo caso $\Gamma$ è elevato.
\item \textbf{Convezione inefficiente:} Il valore di $U$ è grande e $\Gamma$ è piccolo, indicando elevate perdite di calore lungo il cammino.
\item \textbf{Casi intermedi:} È necessario valutare attentamente il bilancio tra l'energia iniziale dell'elemento e la quota dissipata lungo il percorso convettivo.
\end{enumerate}

\section{Applicazioni Astrofisiche della Teoria}

Sebbene l'esposizione si sia concentrata sulle condizioni della \textbf{zona convettiva solare}, i meccanismi descritti dalla \textit{Mixing Length Theory} si applicano a un'ampia varietà di oggetti astrofisici:
\begin{itemize}
\item I \textbf{nuclei delle stelle massicce}.
\item Le fasi terminali delle \textbf{supernovae di Tipo II}, fondamentali per determinare le condizioni di formazione di oggetti peculiari come le \textbf{magnetar} (stelle di neutroni caratterizzate dai campi magnetici più intensi dell'universo).
\end{itemize}

È cruciale comprendere che \textbf{esiste un limite fisico alla quantità di energia che la convezione è in grado di trasportare}.
Se il flusso di energia richiesto supera questo limite, la stella perde la condizione di equilibrio idrostatico, innescando processi esplosivi. Un esempio emblematico di questo fenomeno è rappresentato dai \textbf{pulsi termici} nelle stelle del \textbf{Ramo Asintotico delle Giganti (AGB)}.

\section{Esempio Numerico: Calcolo dei Gradienti Solari}

Per tradurre le equazioni in un'intuizione fisica concreta, inseriamo dei valori tipici per quantificare l'efficienza della convezione nel Sole.
Assumiamo le seguenti ipotesi semplificative:
\begin{itemize}
\item Il plasma si comporta come un \textbf{gas perfetto omogeneo} (assenza di peculiarità nel peso molecolare medio).
\item Il calore specifico adimensionato vale $\frac{c_p}{R} = \frac{5}{2}$.
\item Il parametro di mixing length ($\alpha$) è posto pari a $1$, sicché la lunghezza di mescolamento $L_m$ coincide con l'altezza di scala della pressione $H_p$ ($L_m = H_p$).
\end{itemize}

Valutiamo le condizioni a metà del raggio solare ($r = R_\odot / 2$), assumendo per semplicità che tale raggio racchiuda metà della massa solare ($M_r = M_\odot / 2$). Iniziamo con un primo set di valori di riferimento:
\begin{itemize}
\item Temperatura: $T \sim 10^7$ K
\item Opacità di riferimento: $\kappa \sim 1$ cm$^2$/g
\item Densità di riferimento: $\rho \sim 1$ g/cm$^3$
\end{itemize}
In queste condizioni, il parametro $U$ risulta essere dell'ordine di $10^{-8}$.
Considerando che il parametro $W$ ($\nabla_{rad} - \nabla_{ad}$) ha tipicamente valori compresi tra $1$ e $100$, si ricava che la deviazione del gradiente reale rispetto a quello adiabatico è minima:
\begin{equation}
\nabla - \nabla_{ad} \sim 10^{-5} \div 10^{-4}
\end{equation}
Il sistema è già estremamente vicino all'adiabaticità.

Tuttavia, per approssimare in modo più rigoroso le \textbf{reali condizioni interne del Sole}, è necessario correggere i valori di densità e opacità:
\begin{itemize}
\item Densità reale: $\rho \sim 10^2$ g/cm$^3$
\item Opacità reale: $\kappa \sim 10$ cm$^2$/g
\end{itemize}
Ricalcolando i parametri con questi nuovi dati, il valore di $U$ crolla a $\sim 10^{-13}$. Di conseguenza, la differenza tra i gradienti diviene ancora più trascurabile:
\begin{equation}
\nabla - \nabla_{ad} \sim 10^{-7}
\end{equation}
Possiamo dunque affermare con sicurezza che, \textbf{nel cuore della zona convettiva solare, il gradiente reale è di fatto coincidente con il gradiente adiabatico}. Il parametro $\Gamma$ assume valori enormi, compresi tra $10^6$ e $10^9$, confermando l'estrema efficienza della convezione.

Questa approssimazione \textbf{cade in prossimità delle interfacce} (sia al limite inferiore che superiore della zona convettiva), dove gli effetti di bordo diventano dominanti. In tali regioni, per uno studio accurato, diverrebbe interessante calcolare la \textbf{scala temporale di turnover convettivo} (\textit{convective turnover timescale}), sebbene la sua derivazione analitica risulti complessa e vada oltre gli scopi di questa trattazione.

\section{Il Regime Convettivo nel Sole}

Per avere un'idea quantitativa, pur approssimata, dell'entità dei fenomeni, consideriamo la differenza tra il gradiente termico dell'elemento convettivo e quello dell'ambiente circostante. Nelle condizioni della zona convettiva solare, questa differenza è estremamente piccola, dell'ordine di:
\begin{equation}
\Delta \nabla \approx 10^{-12} \div 10^{-10} \text{ K/cm}
\end{equation}
Questo si traduce in una differenza di temperatura effettiva tra la cella convettiva e l'ambiente circostante tipicamente intorno ai $10^{-21}$ K, un valore microscopico.

Nonostante questa differenza di temperatura infinitesimale, la \textbf{velocità dell'elemento convettivo} risulta essere dell'ordine di:
\begin{equation}
v \approx 1 \div 100 \text{ m/s}
\end{equation}
Con questa velocità, il \textbf{tempo di vita} (o \textit{lifetime}) dell'elemento in risalita si attesta tra 1 e 100 giorni.

Per avere un termine di paragone, la \textbf{velocità del suono} nel mezzo in quelle condizioni è molto più alta, sicché la velocità convettiva è circa $10^{-6} \div 10^{-4}$ volte la velocità del suono locale. Il moto è pertanto \textbf{fortemente subsonico}.

In conclusione, nelle condizioni fisiche del Sole, \textbf{differenze di temperatura infinitesimali possono dare origine a un trasporto di energia estremamente efficiente}. Non stiamo parlando di differenze di 100 K tra la cella e l'ambiente, ma di frazioni minime di grado. L'elemento convettivo può mantenere questo piccolissimo dislivello termico per giorni, risalendo lentamente a pochi metri al secondo, e ciò è sufficiente per garantire una convezione altamente efficace.

\subsection{Tipologie di Convezione: Turbolenta vs Laminare}

Per completare il quadro sulla convezione, dobbiamo considerare la natura del moto fluido che si instaura nel plasma. Bisogna ricordare che il modello della \textbf{Teoria della Lunghezza di Mescolamento (MLT)} appena descritto, seppur formalizzato in un sistema di cinque equazioni e molteplici variabili, costituisce un quadro \textbf{fortemente semplificato}.
L'assunto di base - un elemento che risale trasportando energia per una lunghezza caratteristica determinata - è una schematizzazione estrema della reale complessità idrodinamica.

Tuttavia, tale semplificazione è straordinariamente efficace per lo studio dell'\textbf{evoluzione stellare}, in quanto permette di modellare oggetti radicalmente diversi: dalle nane rosse alle supergiganti, dalle stelle fredde a quelle caldissime, riuscendo di fatto a descrivere l'intera casistica del \textbf{Diagramma H-R} (Hertzsprung-Russell).

Da un punto di vista puramente idrodinamico, la convezione assume forme diverse a seconda del mezzo:
\begin{itemize}
\item Il plasma stellare parzialmente o totalmente ionizzato (come nel Sole).
\item I fluidi liquidi, come l'acqua in ebollizione.
\item L'interno dei pianeti rocciosi, caratterizzato da moti convettivi nel magma del mantello.
\end{itemize}
Benché discipline come la \textbf{Magnetoidrodinamica (MHD)} offrano strumenti per descrivere questi scenari con maggiore rigore, per gli scopi dell'evoluzione stellare ci limiteremo a classificare la convezione in due regimi principali, definiti attraverso il \textbf{Numero di Reynolds} ($Re$):
\begin{equation}
Re = \frac{v \cdot \rho \cdot L_m}{\eta}
\end{equation}
dove:
\begin{itemize}
\item $v$ è la velocità dell'elemento
\item $\rho$ è la densità
\item $L_m$ è la mixing length (lunghezza caratteristica)
\item $\eta$ è la viscosità dinamica del mezzo
\end{itemize}

I due regimi sono i seguenti:
\begin{enumerate}
\item \textbf{Convezione Turbolenta:} Nel plasma solare e stellare in generale, la viscosità è relativamente bassa, portando a Numeri di Reynolds enormi ($Re \gg 1$).
\item \textbf{Convezione Laminare:} Nel magma planetario, la viscosità ($\eta$) e la densità ($\rho$) sono altissime, e le velocità bassissime. Questo si traduce in Numeri di Reynolds molto piccoli.
\end{enumerate}

\section{Stratificazione Termica e Struttura Stellare}

A prescindere dalla complessità del regime idrodinamico, i codici numerici per il calcolo della struttura stellare utilizzano le equazioni del trasporto energetico (risolvendo l'equazione cubica vista in precedenza) per derivare il \textbf{gradiente reale} ($\nabla$). Tale gradiente definisce la struttura stellare in ogni dato momento dell'evoluzione, descrivendo la variazione della temperatura in funzione della frazione di massa o del raggio.

Per comprendere qualitativamente i risultati di questa stratificazione, si può considerare l'equazione dell'\textbf{equilibrio idrostatico} - ricordando che l'instabilità convettiva non viola globalmente l'equilibrio idrostatico della stella.
La variazione di pressione tra due layer contigui in equilibrio idrostatico è espressa da:
\begin{equation}
dP = - \frac{G M_r}{r^2} \rho , dr
\end{equation}
Questa espressione può essere riscritta passando a coordinate sferiche inverse:
\begin{equation}
dP = G M_r \rho , d\left(\frac{1}{r}\right)
\end{equation}
Poiché il gradiente di temperatura $\nabla$ è definito come la derivata logaritmica rispetto alla pressione ($\nabla \equiv \frac{d \ln T}{d \ln P}$), il differenziale esatto di temperatura $dT$ si ottiene come:
\begin{equation}
dT = \nabla , T , \frac{dP}{P}
\end{equation}
Assumendo l'Equazione di Stato dei \textbf{gas perfetti}, sappiamo che:
\begin{equation}
\frac{T \rho}{P} = \frac{\mu}{R_{\text{gas}}}
\end{equation}
dove $\mu$ è il peso molecolare medio e $R_{\text{gas}}$ è la costante dei gas (o costante universale dei gas). Sostituendo l'espressione per $dP$ e l'equazione dei gas perfetti nell'equazione di $dT$, otteniamo:
\begin{equation}
dT = \nabla , \frac{G \mu}{R_{\text{gas}}} M_r , d\left(\frac{1}{r}\right)
\end{equation}
Integrando questa espressione, la differenza di temperatura tra due shell distinte di raggi $r_1$ e $r_2$ (con $r_1 < r_2$) sarà:
\begin{equation}
\Delta T \approx \frac{G M_r \mu}{R_{\text{gas}}} \left( \frac{1}{r_1} - \frac{1}{r_2} \right)
\end{equation}
Questa relazione mostra in prima approssimazione come avviene la stratificazione termica, svelando che la temperatura decresce radialmente in proporzione a $\frac{1}{r}$ all'interno della stella.

\section{Introduzione all'Evoluzione Stellare e alle Reazioni Nucleari}

Con la definizione dei meccanismi di trasporto, abbiamo completato la discussione sulle equazioni basilari che definiscono la \textbf{struttura stellare}. La risoluzione di questo set di equazioni equivale a ottenere un'\textit{istantanea} della stella in un dato istante temporale.

Il fattore fondamentale che guida il \textbf{cambiamento della struttura nel tempo}, ovvero l'evoluzione stellare vera e propria, è la variazione del \textbf{peso molecolare medio ($\mu$)}, indotta dalle \textbf{reazioni nucleari nel nucleo stellare}.
Grazie alla disparità tra le scale temporali (la scala temporale nucleare è ordini di grandezza maggiore rispetto a quella dinamica e termica), l'evoluzione stellare può essere studiata come una \textbf{sequenza di istantanee di equilibrio}. Ad eccezione di eventi su scale temporali estremamente rapide (es. Flash dell'Elio o Supernovae), è sufficiente calcolare il tasso di variazione di $\mu$ nel nucleo e aggiornare passo-passo la struttura. Questo è esattamente il \textit{modus operandi} dei codici di simulazione evolutiva.

\subsection{Sensibilità alla Temperatura delle Catene di Combustione dell'Idrogeno}

L'energia prodotta dalle reazioni nucleari viene calcolata approfonditamente attraverso complessi network di reazioni. Tuttavia, ai fini della comprensione dell'evoluzione, è sufficiente descrivere il \textbf{tasso di produzione di energia nucleare} ($\epsilon$) tramite una \textbf{legge di potenza empirica} dipendente dalla densità ($\rho$) e dalla temperatura ($T$):
\begin{equation}
\epsilon = \epsilon_0 , \rho^{\lambda} , T^{\nu}
\end{equation}
dove $\epsilon_0$ è una costante, $\lambda$ l'esponente della densità e $\nu$ l'esponente della temperatura. Questi esponenti dipendono dal tipo di ciclo nucleare predominante nel nucleo.

Considerando il bruciamento dell'Idrogeno (fase di \textbf{Sequenza Principale}), si presentano due canali principali, caratterizzati da sensibilità estremamente diverse rispetto alla temperatura ($\nu$):

\begin{enumerate}
\item \textbf{Catena Protone-Protone (PP):} Dominante alle temperature inferiori. In questo caso $\lambda \approx 1$ e \textbf{$\nu \approx 4$}. La dipendenza dalla temperatura è significativa ma moderata.


\item \textbf{Ciclo CNO:} Dominante a temperature più elevate. La dipendenza dalla densità resta invariata ($\lambda \approx 1$), ma la sensibilità termica esplode: \textbf{$\nu \approx 15 \div 18$}.



\end{enumerate}

\subsection{Impatto sulla Struttura del Nucleo: Stelle di Piccola e Grande Massa}

La dicotomia nella sensibilità termica ($\nu$) determina un profondo divario strutturale tra stelle di diversa massa sin dal momento dell'ingresso in Sequenza Principale:

\begin{itemize}
\item \textbf{Stelle di Piccola Massa ($M \lesssim 1.4 M_\odot$):}
Le temperature interne più basse favoriscono la \textbf{Catena PP}. Essendo $\nu \approx 4$, la produzione di energia decresce dolcemente procedendo dal centro (più caldo) verso gli strati esterni del nucleo. Il gradiente di flusso energetico non è eccessivamente ripido e può essere interamente trasportato dal trasporto radiativo.
\textbf{Conseguenza:} Le stelle di piccola massa presentano un \textbf{nucleo radiativo}.


\item \textbf{Stelle di Grande Massa ($M \gtrsim 1.4 M_\odot$):} 
Le elevate temperature centrali innescano il \textbf{Ciclo CNO}. Avendo una dipendenza $\nu \approx 18$, una minima variazione radiale di temperatura causa un crollo verticale della produzione di energia. Tutta l'energia viene generata in un volume piccolissimo ed estremamente caldo al centro della stella. Questo genera un differenziale di flusso enorme tra gli strati, per smaltire il quale il solo trasporto radiativo risulta insufficiente, superando le condizioni di stabilità ($\nabla_{rad} > \nabla_{ad}$). 
\textbf{Conseguenza:} Le stelle di grande massa sviluppano un \textbf{nucleo convettivo}.

\end{itemize}

\subsection{Il Bruciamento dell'Elio: Il Processo Tre Alfa ($3\alpha$)}

Per completare questo schema e comprendere appieno le implicazioni fisiche per la struttura stellare, è necessario considerare anche il \textbf{processo tre alfa} ($3\alpha$), ovvero la reazione nucleare responsabile del bruciamento dell'elio.

Analizzando questo processo attraverso la medesima legge di potenza empirica per il tasso di produzione di energia ($\epsilon \propto \rho^{\lambda} \, T^{\nu}$), i parametri assumono i seguenti valori tipici:
\begin{itemize}
\item L'esponente della densità risulta $\lambda \approx 2$.
\item L'esponente della temperatura (erroneamente pronunciato come "Ni" nella trascrizione, ma corrispondente alla lettera greca $\nu$) raggiunge valori estremi: \textbf{$\nu \approx 40$}.
\end{itemize}

La situazione, quando si instaura il regime di bruciamento dell'elio, risulta pertanto ancora più estrema rispetto al Ciclo CNO. La sensibilità alle minime variazioni di temperatura è talmente elevata che la produzione di energia collassa a zero non appena ci si allontana dal centro esatto del nucleo.

Ciò implica che, durante l'innesco e lo svolgimento del processo $3\alpha$, le probabilità che il trasporto radiativo divenga instabile e che si instauri un \textbf{regime di trasporto convettivo dell'energia} sono altissime.

\section{Conclusione}

Con queste considerazioni si chiude l'analisi dettagliata dei meccanismi di trasporto energetico e dei gradienti termici. Avendo ora a disposizione tutte le equazioni fondamentali per descrivere la struttura interna di una stella (le "istantanee" evolutive), dalla prossima lezione si inizierà ad affrontare in modo diretto e sistematico l'evoluzione stellare nel tempo.